% Integrated 4 September 2026 from the exact 10^8 census and independent replay.
% This section retains the literal ternary equation, chamber, strict wall, height,
% and original-coordinate inverse.

\section{Exact unit-orbit census through height \(10^8\)}
\label{koide-count-sec-census-1e8}

The census retains exactly the same equation, chamber, strict wall, and
height as the preceding theorems:
\begin{equation}
 (a,q,n)\in\mathbb Z^3,\qquad
 a^2-3q^2-18n^2=-2,\qquad
 a>0,\quad q,n\geq0,\quad a+4n>q,\qquad
 4a+17n\leq X.
 \label{koide-count-eq-census-domain}
\end{equation}
It does not infer an asymptotic from finite data.  Its exact fixed-\(n\)
equation and factorization are
\begin{equation}
 a^2-3q^2=m_n,\qquad
 m_n=18n^2-2=2(3n-1)(3n+1).
 \label{koide-count-eq-census-fixed-n}
\end{equation}
For each \(n\geq1\), every ideal of norm \(m_n\) is enumerated by its
Hermite-normal-form triple
\begin{equation}
 (C,D,E)\in\mathbb Z_{>0}^2\times\mathbb Z,\qquad
 C^2D=m_n,\qquad
 0\leq E<D,\qquad E^2\equiv3\pmod D.
 \label{koide-count-eq-census-hnf}
\end{equation}
The generator is reduced by exact multiplication by the fundamental unit
\(2+\sqrt3\) to the unique seed \(a_0>2q_0\geq0\).  Its forward unit
orbit is
\begin{equation}
 \binom{a_{k+1}}{q_{k+1}}
 =
 \begin{pmatrix}2&3\\1&2\end{pmatrix}
 \binom{a_k}{q_k},\qquad k\geq0.
 \label{koide-count-eq-census-unit-step}
\end{equation}
Each emitted point is reconstructed by
\begin{equation}
 x=\frac{a+4n-q}{2},\qquad
 y=\frac{a+4n+q}{2},\qquad
 z=4a+17n,
 \label{koide-count-eq-census-inverse}
\end{equation}
and is rechecked in both coordinate systems and in the original polynomial
\(\Delta(x,y,z)=-1\).  Thus the accelerated enumeration does not replace
the original objects by a floating-point or rescaled surrogate.

\subsection{Complete result and checkpoint table}

The complete integer-domain audit for
\(1\leq X\leq100{,}000{,}000\) gives
\begin{equation}
 \boxed{R(100{,}000{,}000)=3{,}536{,}002.}
 \label{koide-count-eq-census-final-count}
\end{equation}
The endpoint in \(n\) is \(2{,}943{,}725\); the enumeration contains
\(2{,}327{,}513\) canonical unit orbits.  The following checkpoints are
exact integer counts.  The residuals use the proved constant
\(C_0\) in \eqref{koide-count-eq-final-constant}.
\begin{center}
\begin{tabular}{r@{\qquad}r@{\qquad}r@{\qquad}r}
\toprule
\(X\)&\(R(X)\)&\(R(X)-C_0X\)&\((R(X)-C_0X)/\sqrt X\)\\
\midrule
200,000&7,109&\(36.590440053629\)&\(0.081818711287\)\\
250,000&8,854&\(13.488050067037\)&\(0.026976100134\)\\
500,000&17,695&\(13.976100134073\)&\(0.019765190359\)\\
1,000,000&35,381&\(18.952200268147\)&\(0.018952200268\)\\
2,000,000&70,637&\(-87.095599463706\)&\(-0.061585888992\)\\
5,000,000&176,943&\(132.761001340734\)&\(0.059372524752\)\\
10,000,000&353,598&\(-22.477997318532\)&\(-0.007108166877\)\\
20,000,000&707,458&\(217.044005362937\)&\(0.048532515010\)\\
50,000,000&1,768,230&\(127.610013407342\)&\(0.018046781166\)\\
100,000,000&3,536,002&\(-202.779973185317\)&\(-0.020277997319\)\\
\bottomrule
\end{tabular}
\end{center}

\subsection{Certified extrema over every integer}

Write \(E(X)=R(X)-C_0X\).  Between consecutive occupied heights,
\(R(X)\) is constant and both \(E(X)\) and
\(E(X)/\sqrt X\) are strictly decreasing, since
\begin{equation}
 \frac{d}{dX}\left(\frac{R}{\sqrt X}-C_0\sqrt X\right)
 =-\frac{R}{2X^{3/2}}-\frac{C_0}{2\sqrt X}<0.
 \label{koide-count-eq-census-monotonicity}
\end{equation}
Checking both integer endpoints of every constant-count interval therefore
certifies extrema over every integer in the audited range, rather than at
selected checkpoints.  On \(1\leq X\leq10^8\), the maximum signed residual
is
\begin{equation}
 \begin{aligned}
 E(77{,}078{,}404)
 &=721.7934249571297569656681553538691093379\ldots,\\
 R(77{,}078{,}404)&=2{,}726{,}372,
 \end{aligned}
 \label{koide-count-eq-census-max-positive}
\end{equation}
and the minimum signed, hence largest absolute, residual is
\begin{equation}
 \begin{aligned}
 E(32{,}461{,}294)
 &=-812.8300691488068897091966006797491655933\ldots,\\
 R(32{,}461{,}294)&=1{,}147{,}085.
 \end{aligned}
 \label{koide-count-eq-census-max-negative}
\end{equation}
The maximum of
\(\lvert E(X)\rvert/\sqrt X\) over the full audited range is attained at
\(X=160\), where
\[
 \frac{\lvert E(160)\rvert}{\sqrt{160}}
 =0.185157076934082141384010185515317233816\ldots.
\]
On \(10{,}000\leq X\leq10^8\), the maximum is
\begin{equation}
 \begin{aligned}
 \frac{\lvert E(867{,}240)\rvert}{\sqrt{867{,}240}}
 &=0.171400192808814559728949658032345552502\ldots,\\
 R(867{,}240)&=30{,}827,
 \end{aligned}
 \label{koide-count-eq-census-tail-extremum}
\end{equation}
and the root entering at that height is
\begin{equation}
 (a,q,n)=(199997,115061,3956),\qquad
 (x,y,z)=(50380,165441,867240).
 \label{koide-count-eq-census-tail-witness}
\end{equation}
The bound \(\lvert E(X)\rvert\leq0.171\sqrt X\) first fails on this tail at
\(X=867{,}216\), where \(R(X)=30{,}826\) and
\[
 E(X)=159.4663553077412220330449353234319817199\ldots.
\]
No counterexample to the finite statement with constant \(0.172\) occurs
through \(10^8\).  These are bounded computations, not a claim that either
constant persists.

The audit finds \(12{,}916\) sign changes.  The \(3{,}536{,}002\) roots
occupy \(2{,}472{,}015\) distinct heights; the maximum shell multiplicity
is \(35\), at \(z=92{,}923{,}120\), and the largest gap between consecutive
occupied heights is \(564\), from \(71{,}586{,}276\) to \(71{,}586{,}840\).

\subsection{Independent replay and exact algorithm}

An independent pass decompressed and rechecked all \(3{,}536{,}002\)
six-coordinate rows, recomputed every checkpoint, verified the coordinate
maps, equations, ordering, chamber, height, nonduplication, and gzip CRC,
and returned \texttt{PASS}.  It also compared the literal \((n,a)\) scan
through \(X=20{,}000\), every generalized-Pell seed for \(1\leq n\leq1000\),
and the full six-coordinate point set through \(X=200{,}000\).

The reproducibility sources and receipts are retained project-relatively under
\path{supporting_materials/workbench/evidence/koide_exact_checkpoint_20260904/census}.
The two \(10^8\)-range outputs and the full coordinate receipt are in its
\path{final_100000000} subdirectory.  The principal digests are
\begin{center}
\begin{tabular}{>{\raggedright\arraybackslash}p{0.43\linewidth}
                >{\raggedright\arraybackslash}p{0.51\linewidth}}
\toprule
artifact&SHA-256\\
\midrule
\texttt{koide\_fast\_enum.py}&
\texttt{dddb73b828b63be08e2ff0b5346edb32\allowbreak%
31662005f83d5b2448c06fef6ba31017}\\
\texttt{koide\_discrepancy\_100000000.json}&
\texttt{da92f0866074e530b8174a48915caa0f\allowbreak%
65cd4cd8957e209306adc200739d1961}\\
\texttt{koide\_roots\_through\_\allowbreak 100000000.csv.gz}&
\texttt{dfc313492aedb396dcfa10520bfefc699\allowbreak%
ba601357bc9933d921dfa3884310bd4}\\
\texttt{coordinate\_audit\_receipt.json}&
\texttt{3010363108e7b460b7331a92f3263bbd\allowbreak%
77992aedfc3064973750e16d707da80f}\\
\bottomrule
\end{tabular}
\end{center}

The exact algorithm first finds the largest admissible \(n\) by integer
binary search, then factors the unchanged integers \(3n-1\) and \(3n+1\)
with a smallest-prime-factor sieve.  It enumerates every HNF triple in
\eqref{koide-count-eq-census-hnf}, constructs all roots of \(3\) modulo
\(D\) by prime-power lifting and the Chinese remainder theorem, reduces
each generator to the unique seed, and iterates
\eqref{koide-count-eq-census-unit-step} until the literal height fails.
For each \(n\), the independently checked identity
\[
 \#\mathcal P(m_n)=\rho(m_n)=\rho(3n-1)\rho(3n+1)
\]
certifies the orbit multiplicity.  The work is
\[
 X^{1+o(1)}+O(R(X))
\]
integer operations before the optional height sort.  This complexity and
the finite census use no conjectural discrepancy estimate.
